% UC-071: Concept part (CWW assistant design).
\section{Concept: A Computing-with-Words Assistant}
\label{sec:concept}
\placeholder{UC-071: vocabulary, truth of summaries, quality measures, decision rule and
faithful explanation (docs/methods.md sections 1--6).}

\subsection{Vocabulary and Truth}
\placeholder{UC-071.}

\subsection{Decision and Explanation}
\placeholder{UC-071.}


\subsection{Methodological Baseline}
\label{sec:methodological-baseline}
We apply the unrestricted linguistic-summary truth computation of
\citet[Eq.~2]{pincay2021} as a methodological baseline for Swiss train arrival delays.
The input is the study's measured Ist-Daten arrivals, without a passenger-specific
restriction, and the summarizers are the study's four arrival-delay words: on time,
slightly late, late and very late. For each word, its fuzzy proportion is evaluated
with the quantifiers few, about half and most. We compute these scores for the pooled
network and separately for each study half.

The baseline uses the dedicated \texttt{pincay2021} vocabulary configuration, whose
curves follow \citet[Fig.~2, p.~172]{pincay2021}. The main study retains its own
quantifier breakpoints (D-028). A summary is valid when its truth is at least the
configured mining threshold $\beta$; equality is included. The default follows
\citet[Section~\citednum{3.2}, p.~172]{pincay2021}, while configuration overrides also apply to
this baseline. This is a methodological comparison on Swiss arrival data, not a
reproduction or numerical comparison of the paper's traffic-message results.
